% ================= 参考文献 =================

\section*{参考文献}\label{sec:references}
\addcontentsline{toc}{section}{参考文献}

\begin{thebibliography}{99}

\bibitem{ref:1} Y.~AMICE \& J.~V\'ELU - Distributions p-adiques associ\'ees aux s\'eries de Hecke, Ast\'erisque 24-25 (1975), p.~119-131.

\bibitem{ref:2} F.~BALDASSARRI \& B.~CHIARELLOTTO - Algebraic versus rigid cohomology with logarithmic coefficients, in \emph{Barsotti Symposium in Algebraic Geometry} (Abano Terme, 1991), Perspect.\ Math., vol.~15, Academic Press, 1994, p.~11-50.

\bibitem{ref:3} L.~BERGER \& C.~BREUIL - Sur quelques repr\'esentations potentiellement cristallines de $\GL_2(\Qp)$, Ast\'erisque 330 (2010), p.~155-211.

\bibitem{ref:4} L.~BERGER \& C.~BREUIL - Towards a p-adic Langlands programme, cours au C.M.S.\ de Hangzhou (2004), http://www.ihes.fr/\textasciitilde breuil/publications.html.

\bibitem{ref:5} P.~BERTHELOT - Cohomologie rigide et cohomologie rigide \`a support propre. Premi\`ere partie, http://www.maths.univ-rennes1.fr/\textasciitilde berthelo/.

\bibitem{ref:6} N.~BOURBAKI - \emph{Topological vector spaces. Chapters 1-5}, Elements of Mathematics (Berlin), Springer, 1987.

\bibitem{ref:7} C.~BREUIL - Sur quelques repr\'esentations modulaires et p-adiques de $\GL_2(\Qp)$. II, J.\ Inst.\ Math.\ Jussieu 2 (2003), p.~23-58.

\bibitem{ref:8} C.~BREUIL - Invariant $L$ et s\'erie sp\'eciale p-adique, Ann.\ Sci.\ Ecole Norm.\ Sup.\ 37 (2004), p.~559-610.

\bibitem{ref:9} C.~BREUIL - S\'erie sp\'eciale p-adique et cohomologie \'etale compl\'et\'ee, ce volume.

\bibitem{ref:10} K.~BUZZARD \& F.~CALEGARI - Slopes of overconvergent 2-adic modular forms, Compos.\ Math.\ 141 (2005), p.~591-604.

\bibitem{ref:11} K.~BUZZARD \& R.~TAYLOR - Companion forms and weight one forms, Ann.\ of Math.\ 149 (1999), p.~905-919.

\bibitem{ref:12} R.~F.~COLEMAN - A p-adic Shimura isomorphism and p-adic periods of modular forms, in \emph{p-adic monodromy and the Birch and Swinnerton-Dyer conjecture} (Boston, MA, 1991), Contemp.\ Math., vol.~165, Amer.\ Math.\ Soc., 1994, p.~21-51.

\bibitem{ref:13} R.~F.~COLEMAN - Classical and overconvergent modular forms, Invent.\ Math.\ 124 (1996), p.~215-241.

\bibitem{ref:14} R.~F.~COLEMAN - Classical and overconvergent modular forms of higher level, J.\ Th\'eor.\ Nombres Bordeaux 9 (1997), p.~395-403.

\bibitem{ref:15} R.~F.~COLEMAN - p-adic Banach spaces and families of modular forms, Invent.\ Math.\ 127 (1997), p.~417-479.

\bibitem{ref:16} R.~F.~COLEMAN \& A.~IOVITA - Hidden structures on semi-stable curves, ce volume.

\bibitem{ref:17} R.~F.~COLEMAN \& B.~MAZUR - The eigencurve, in \emph{Galois representations in arithmetic algebraic geometry} (Durham, 1996), London Math.\ Soc.\ Lecture Note Ser., vol.~254, Cambridge Univ.\ Press, 1998, p.~1-113.

\bibitem{ref:18} P.~COLMEZ - Une correspondance de Langlands p-adique pour les repr\'esentations semi-stables de dimension 2, pr\'epublication, 2004.

\bibitem{ref:19} P.~COLMEZ - La s\'erie principale unitaire de $\GL_2(\Qp)$, Ast\'erisque 330 (2010), p.~213-262.

\bibitem{ref:20} P.~DELIGNE - Formes modulaires et repr\'esentations de $\GL(2)$, in \emph{Modular functions of one variable, II} (Proc.\ Internat.\ Summer School, Univ.\ Antwerp, Antwerp, 1972), Springer, 1973, p.~55-105. Lecture Notes in Math., Vol.~349.

\bibitem{ref:21} M.~EMERTON - p-adic $L$-functions and unitary completions of representations of p-adic reductive groups, Duke Math.\ J.\ 130 (2005), p.~353-392.

\bibitem{ref:22} M.~EMERTON - Jacquet modules of locally analytic representations of p-adic reductive groups. I.\ Construction and first properties, Ann.\ Sci.\ Ecole Norm.\ Sup.\ 39 (2006), p.~775-839.

\bibitem{ref:23} M.~EMERTON - A local-global compatibility conjecture in the p-adic Langlands programme for $\GL_2/\Q$, Pure Appl.\ Math.\ Q.\ 2 (2006), p.~279-393.

\bibitem{ref:24} M.~EMERTON - On the interpolation of systems of eigenvalues attached to automorphic Hecke eigenforms, Invent.\ Math.\ 164 (2006), p.~1-84.

\bibitem{ref:25} M.~EMERTON - Jacquet modules of locally analytic representations of p-adic reductive groups II.\ The relation to parabolic induction, \`a para\^itre dans J.\ Institut Math.\ Jussieu.

\bibitem{ref:26} M.~EMERTON - Locally analytic vectors in representations of locally p-adic analytic groups, \`a para\^itre dans Memoirs of the AMS.

\bibitem{ref:27} G.~FALTINGS - Crystalline cohomology of semistable curve---the $\Qp$-theory, J.\ Algebraic Geom.\ 6 (1997), p.~1-18.

\bibitem{ref:28} G.~FALTINGS - Almost \'etale extensions, Ast\'erisque 279 (2002), p.~185-270.

\bibitem{ref:29} G.~FALTINGS \& B.~W.~JORDAN - Crystalline cohomology and $\GL(2,\Q)$, Israel J.\ Math.\ 90 (1995), p.~1-66.

\bibitem{ref:30} E.~GHATE - Ordinary forms and their local Galois representations, \`a para\^itre dans les Actes de la conf\'erence \`a Hyderabad 2003.

\bibitem{ref:31} E.~GHATE \& A.~M\'EZARD - Filtered modules attached to modular forms, pr\'epublication, 2006.

\bibitem{ref:32} F.~Q.~GOUV\^EA - \emph{Arithmetic of p-adic modular forms}, Lecture Notes in Math., vol.~1304, Springer, 1988.

\bibitem{ref:33} B.~H.~GROSS - A tameness criterion for Galois representations associated to modular forms (mod $p$), Duke Math.\ J.\ 61 (1990), p.~445-517.

\bibitem{ref:34} H.~HIDA - Galois representations into $\GL_2(\Zp[[X]])$ attached to ordinary cusp forms, Invent.\ Math.\ 85 (1986), p.~545-613.

\bibitem{ref:35} N.~M.~KATZ - p-adic properties of modular schemes and modular forms, in \emph{Modular functions of one variable, III} (Proc.\ Internat.\ Summer School, Univ.\ Antwerp, Antwerp, 1972), Springer, 1973, p.~69-190. Lecture Notes in Mathematics, Vol.~350.

\bibitem{ref:36} N.~M.~KATZ \& B.~MAZUR - \emph{Arithmetic moduli of elliptic curves}, Annals of Math.\ Studies, vol.~108, Princeton Univ.\ Press, 1985.

\bibitem{ref:37} M.~KURIHARA - A tameness criterion, in \emph{Volume \`a la m\'emoire d'O.\ Hyodo}, universit\'e de Tokyo, 1993, p.~27-36.

\bibitem{ref:38} B.~MAZUR - Modular curves and the Eisenstein ideal, Publ.\ Math.\ I.H.\'E.S.\ 47 (1977), p.~33-186.

\bibitem{ref:39} K.~A.~RIBET - Galois representations attached to eigenforms with Nebentypus, in \emph{Modular functions of one variable, V} (Proc.\ Second Internat.\ Conf., Univ.\ Bonn, Bonn, 1976), Springer, 1977, p.~17-51. Lecture Notes in Math., Vol.~601.

\bibitem{ref:40} T.~SAITO - Modular forms and p-adic Hodge theory, Invent.\ Math.\ 129 (1997), p.~607-620.

\bibitem{ref:41} W.~H.~SCHIKHOF - \emph{Non-Archimedean calculus}, Lecture notes, vol.~7812, Katholieke Universiteit Mathematisch Instituut, 1978.

\bibitem{ref:42} W.~H.~SCHIKHOF - \emph{Ultrametric calculus}, Cambridge Studies in Advanced Math., vol.~4, Cambridge Univ.\ Press, 1984.

\bibitem{ref:43} P.~SCHNEIDER \& J.~TEITELBAUM - Banach space representations and Iwasawa theory, Israel J.\ Math.\ 127 (2002), p.~359-380.

\bibitem{ref:44} P.~SCHNEIDER \& J.~TEITELBAUM - Locally analytic distributions and p-adic representation theory, with applications to $\GL_2$, J.\ Amer.\ Math.\ Soc.\ 15 (2002), p.~443-468.

\bibitem{ref:45} P.~SCHNEIDER \& J.~TEITELBAUM - Algebras of p-adic distributions and admissible representations, Invent.\ Math.\ 153 (2003), p.~145-196.

\bibitem{ref:46} J-P.~SERRE - Quelques applications du th\'eor\`eme de \v{C}ebotarev, Publ.\ Math.\ I.H.E.S.\ 54 (1981), p.~123-201.

\bibitem{ref:47} M.~VISHIK - Nonarchimedean measures connected with Dirichlet series, Math.\ USSR Sb.\ 28 (1976), p.~216-228.

\end{thebibliography}

\vspace{1.5em}

\noindent C.~BREUIL, D\'epartement de Math\'ematiques, B\^atiment 425, Facult\'e des Sciences d'Orsay, Universit\'e Paris-Sud 11, 91405 Orsay Cedex, France \quad$\bullet$\quad E-mail: breuil@ihes.fr

\noindent M.~EMERTON, Mathematics Department, Northwestern University, 2033 Sheridan Road, Evanston, IL 60208, U.S.A. \quad$\bullet$\quad E-mail: emerton@math.northwestern.edu
